\documentclass[12pt,a4paper]{article}
\usepackage[utf8]{inputenc}
\usepackage{graphicx}
\usepackage{geometry}
\usepackage{hyperref}
\geometry{a4paper, margin=2.5cm}

\title{\textbf{Tugas Praktikum: Website Katalog Sederhana dengan Django}}
\author{
    \textbf{Faris Zain As-Shadiq} \\
    NPM: 2308107010039
}
\date{\today}

\begin{document}

\maketitle

\section*{Detail Pengumpulan}
\begin{itemize}
    \item \textbf{Nama:} Faris Zain As-Shadiq
    \item \textbf{NPM:} 2308107010039
    \item \textbf{Link GitHub:} \href{https://github.com/fariszain/Katalog_Faris.git}{https://github.com/fariszain/Katalog\_Faris.git}
\end{itemize}

\section*{Penjelasan Singkat Program}
Proyek ini adalah sebuah website katalog produk sederhana yang dibangun menggunakan framework Django. Website ini bertujuan untuk menampilkan daftar produk beserta detailnya. Data produk yang ditampilkan menggunakan data \textit{hardcoded} di \texttt{views.py} (tidak menggunakan database). Terdapat empat halaman utama yaitu Homepage, Daftar Produk, Detail Produk, dan Kontak. Seluruh mekanisme \textit{routing} memanfaatkan sistem \texttt{urls.py} dan \textit{template} HTML dari Django.

\newpage
\section*{Lampiran Screenshot}

\subsection*{1. Homepage (URL: \texttt{/})}
\begin{figure}[h!]
    \centering
    % Pastikan Anda mengganti 'home.png' dengan nama file screenshot Anda
    \includegraphics[width=0.8\textwidth]{home.png}
    \caption{Tampilan Homepage}
\end{figure}

\subsection*{2. Halaman Daftar Produk (URL: \texttt{/produk/})}
\begin{figure}[h!]
    \centering
    % Pastikan Anda mengganti 'produk.png' dengan nama file screenshot Anda
    \includegraphics[width=0.8\textwidth]{produk.png}
    \caption{Tampilan Daftar Produk (Menampilkan 3 Produk)}
\end{figure}

\subsection*{3. Halaman Detail Produk (URL: \texttt{/produk/<id>/})}
\begin{figure}[h!]
    \centering
    % Pastikan Anda mengganti 'detail.png' dengan nama file screenshot Anda
    \includegraphics[width=0.8\textwidth]{detail.png}
    \caption{Tampilan Detail Produk Berdasarkan ID}
\end{figure}

\subsection*{4. Halaman Kontak (URL: \texttt{/kontak/})}
\begin{figure}[h!]
    \centering
    % Pastikan Anda mengganti 'kontak.png' dengan nama file screenshot Anda
    \includegraphics[width=0.8\textwidth]{kontak.png}
    \caption{Tampilan Halaman Kontak}
\end{figure}

\end{document}
